\section{为什么 recurrent / linear models 又回来了}

这堂课表面上讨论 State Space Model（SSM，状态空间模型）与 Transformer 的效率取舍，真正要回答的却是更基础的问题：\textbf{一个自回归模型用什么形式保存过去？} 如果过去被保存为逐 token 可寻址的缓存，模型擅长精确查找；如果过去被压缩进固定状态，模型更适合在线更新与抽象。二者的复杂度差异只是这种记忆接口的外在结果，而不是全部本质。

\subsection{从研究潮流到生产系统}

先看讲者为什么把这件事称为“复兴”。Mamba 之后出现了 xLSTM、DeltaNet、Gated DeltaNet、Test-Time Training 等多条路线；更重要的是，它们已经进入 Jamba、Zamba、Samba、Hunyuan、Qwen、Kimi Linear、OLMo Hybrid、Megatron--NeMo 等大型混合模型。读图时不要急着比较名字，而要先辨认两层信息：左边是机制家族，右边是这些机制如何进入 production-scale hybrid model。

\lecturefigure{slide-002.jpg}{现代 recurrent / linear model 家族及其在大型混合模型中的采用}{Stanford 官方录像，00:03:20--00:04:09}

\teachervoice{00:01:28--00:03:55，Albert Gu 强调这些方法已经广泛进入大规模模型；本课因此不逐篇介绍论文，而是上升到共同特征、共同瓶颈与 Transformer 的根本差异。}

同一个机制在不同社区里会被叫作 SSM、linear attention、linear RNN、modern recurrent model 或 linear model。下图把这些词放在同一页，不是说所有公式完全相同，而是说明本课只关心它们共享的\textbf{有限 recurrent state + 线性/次二次序列计算}接口。所谓 hybrid model，则是在网络深度方向交错放置这类层与 quadratic attention 层。

\lecturefigure{slide-003.jpg}{本课采用的 umbrella taxonomy：SSM、linear attention 与 modern recurrent model}{Stanford 官方录像，00:04:10--00:05:09}

\begin{knowledgebox}{术语消化：名字不同，比较轴相同}
\begin{tabularx}{\textwidth}{L{0.20\textwidth}X X}
\toprule
术语 & 常见出发点 & 本课关心的共同接口 \\
\midrule
SSM & 连续/离散状态空间、结构化转移 & 用有限状态递推压缩历史 \\
Linear attention & 把 attention 核写成可递推统计量 & 不保存完整 $L\times L$ attention matrix \\
Linear RNN & 线性 recurrent update 与并行训练 & 状态更新可组合、可 scan \\
Modern recurrent model & 相对 LSTM/GRU 的新一代 recurrence & 扩大状态、增强 gating、适配加速器 \\
Hybrid model & 工程上交错不同 layer type & recurrent 层负责大部分处理，attention 层补精确检索 \\
\bottomrule
\end{tabularx}
\end{knowledgebox}

\begin{warningbox}{不要把 umbrella label 当作数学等价}
Mamba、Gated DeltaNet、xLSTM 与 TTT 的状态形状、更新式、训练算法和 kernel 都不相同。本课把它们分在一组，只是为了比较“有限压缩状态”与“token-resolution cache”这条主轴；任何具体实现结论仍需回到对应论文和硬件。
\end{warningbox}

\subsection{为什么用 autoregressive modeling 做显微镜}

自回归生成把联合分布分解为逐步条件概率：训练时并行预测下一 token，推理时却必须把刚生成的 token 重新送回模型。因此，解码过程中“步骤之间留下什么”会暴露架构的真实记忆合同。

\[
p(x_{1:L})=\prod_{t=1}^{L}p(x_t\mid x_{<t}),
\qquad
x_{t+1}\sim p_\theta(\cdot\mid x_{\le t}).
\]
其中，$L$ 是序列长度，$x_t$ 是第 $t$ 个 token，$x_{<t}$ 是此前上下文，$p_\theta$ 是参数为 $\theta$ 的模型。训练可以一次处理整段序列；生成必须按 $t=1,2,\ldots$ 顺序推进。

\lecturefigure{slide-004.jpg}{Autoregressive modeling：训练预测下一 token，推理逐步采样并回灌}{Stanford 官方录像，00:05:10--00:06:09}

\teachervoice{00:05:10--00:05:49，讲者选择 autoregressive modeling 有两层原因：它既是当前语言模型的核心范式，也最容易把不同 sequence model 的状态与复杂度暴露出来。}

\begin{importantbox}{本课的比较单位不是 layer，而是 generation state}
比较两个架构时，先问：生成完第 $t$ 步后，模型必须保留什么才能生成第 $t+1$ 步？这个跨步持久对象就是 autoregressive state。后续所有关于 recall、compression、online processing 与 token resolution 的讨论，都从这里展开。
\end{importantbox}

\subsection{本章小结}

现代 recurrent / linear model 的价值已不止是论文中的复杂度改写，而是生产模型的真实设计空间。为了穿过模型名称迷雾，本课固定一个比较镜头：观察自回归生成中跨步保存的状态。下一章将看到，Transformer 把过去保存在逐 token cache 中，SSM 则把过去压进有限 recurrent state；这两个接口同时决定能力、归纳偏置与成本。

